\documentclass[11pt,a4paper]{article}
\usepackage{fontspec,amsmath,amssymb,amsthm,microtype}
\setmainfont{lmroman10-regular.otf}[BoldFont=lmroman10-bold.otf,ItalicFont=lmroman10-italic.otf,BoldItalicFont=lmroman10-bolditalic.otf]
\setmonofont{lmmono10-regular.otf}[Scale=MatchLowercase]
\usepackage[margin=25mm]{geometry}
\usepackage[unicode,hidelinks]{hyperref}
\setlength{\emergencystretch}{3em}
\setlength{\parskip}{0.4em}
\setlength{\parindent}{0pt}
\allowdisplaybreaks[2]
\raggedbottom
\title{\texorpdfstring{The original OBABO position laws satisfy the full sampling bound}{The original OBABO position laws satisfy the full sampling bound}}
\author{LangevinLean mathematical exposition}
\date{}
\begin{document}\maketitle
\noindent Source module: \texttt{Kinetic\allowbreak{}Langevin/\allowbreak{}OBABOPosition.\allowbreak{}lean}. The accompanying HTML pairs this mathematical exposition with the exact Lean statements and records the scope of the formalization.


The actual initial phase law is standard Gaussian velocity independent of the specified position law. Coupling these velocities identically costs nothing. The explicit inverse weighted coordinate map then projects the proved full Markov theorem to the actual position laws, with factor \ensuremath{\surd}2 and finite output transport established separately.


\section{\texorpdfstring{Project the contracting geometry to the physical position}{Project the contracting geometry to the physical position}}


The projection is the position coordinate of the explicitly verified inverse transform. The metric lower bound gives \(\|x-y\|^2\le2\|T(v,x)-T(w,y)\|^2\). Both sides are nonnegative, so the square comparison yields the global \ensuremath{\surd}2 Lipschitz coefficient. Composing this projection with the original weighted transform is exactly the ordinary position projection; this is proved from the inverse identity, not supplied as a law equivalence.



{\raggedright\noindent Lean declarations: \texttt{obabo\allowbreak{}Position\allowbreak{}Phase\_\allowbreak{}measurable}, \texttt{obabo\allowbreak{}Position\allowbreak{}Phase\_\allowbreak{}borel}, \texttt{obabo\allowbreak{}Position\allowbreak{}Projection}, \texttt{obabo\allowbreak{}Position\allowbreak{}Projection\_\allowbreak{}lipschitz}, \texttt{obabo\allowbreak{}Position\allowbreak{}Projection\_\allowbreak{}transform}.\par}

\section{\texorpdfstring{Adjoin the same Gaussian velocity to every initial position coupling}{Adjoin the same Gaussian velocity to every initial position coupling}}


For every coupling of two position laws, append one common independent Gaussian velocity. Their phase difference has zero velocity, so its weighted energy is exactly the squared original position distance. The cross-space random-map theorem takes the actual coupling infimum and yields \(W_{2,E}(\mathcal N\otimes\mu,\mathcal N\otimes\nu)\le W_2(\mu,\nu)\). Product-measure swapping proves that the actual pushforward initialization is the stated independent Gaussian-times-position law. No optimal position coupling is assumed.



{\raggedright\noindent Lean declarations: \texttt{gaussian\allowbreak{}Velocity\_\allowbreak{}weighted\_\allowbreak{}wasserstein2}.\par}

\section{\texorpdfstring{Obtain the unconditional actual position estimate and expose finite transport}{Obtain the unconditional actual position estimate and expose finite transport}}


Start from \(\mathcal N(0,I_d)\otimes\mu\), use the actual independent Gaussian-pair OBABO transition at every step, and let \(\nu_n\) be its actual position marginal. Under the original strong-convexity, smoothness, friction and step-size restrictions, the theorem proves both finite \(W_2(\nu_n,\pi)\) and

\[W_2(\nu_n,\pi)\le\sqrt2\,r^nW_2(\mu,\pi)+\frac{6\sqrt2\gamma}{m}h\sqrt d\left(4\sqrt M+\frac{Mh}{2}\right),\qquad r=\sqrt{1-\frac{mh}{3\gamma}}.\]

The initial weighted distance is bounded by the actual initial position distance using the preceding shared-velocity construction. The full weighted theorem and its finite-output counterpart apply to the implemented Markov laws. The \ensuremath{\surd}2 Lipschitz transport theorem projects both actual output measures; the target projection is the already proved normalized Gibbs marginal. Output finiteness is proved before the real-valued inequality is used, so this endpoint can safely imply RMS consequences. The finite-dimensional index type is arbitrary, allowing direct reuse for the scaled particle configuration.



{\raggedright\noindent Lean declarations: \texttt{obabo\_\allowbreak{}position\_\allowbreak{}sampling\_\allowbreak{}bound}.\par}

\end{document}
